\Introduction

Выпускная квалификационная работа посвящена разработке корпоративной системы управления обучением, предназначенной для организации учебного контента, назначения курсов, контроля результатов и проведения адаптивного тестирования в рамках единой клиент--серверной платформы\cite{metodichka_vkr_ikt_2021}.

Актуальность темы обусловлена тем, что в корпоративной среде обучение сотрудников всё чаще строится как непрерывный процесс, включающий адаптацию новых сотрудников (онбординг), обязательные программы, проверку знаний и повторение проблемных тем. При использовании разрозненных инструментов усложняется управление пользователями и ролями, затрудняется назначение курсов, отсутствует единый контур аналитики и слабо поддерживается персонализация учебной траектории. Дополнительным фактором является распространение мобильного обучения, при котором слушатели проходят курсы короткими сессиями со смартфона и ожидают простой навигации, быстрого доступа к материалам и прозрачной обратной связи по результатам тестирования\cite{mobile_learning_solutions_2024,opigno_mobile_lms_2024}.

Для корпоративных сценариев также существенна возможность обслуживания нескольких организаций в рамках одной платформы с логическим разделением данных, ролей и учебного контента. Такой подход позволяет использовать единое программное решение для нескольких арендаторов, сохраняя изоляцию данных и независимость администрирования\cite{moodle_multitenancy_2024,yojji_multitenant_lms_2025}.

Цель работы --- разработать корпоративную LMS с веб-панелью администрирования и мобильным клиентом слушателя, обеспечивающую управление учебным контентом, назначениями, адаптивным тестированием и аналитикой результатов обучения. Система адресована ролям системного администратора платформы, администратора организации, преподавателя и слушателя.

Для достижения поставленной цели необходимо решить следующие задачи:
\begin{parenlist}
  \item проанализировать предметную область корпоративного обучения и исходное состояние процессов;
  \item выявить основные проблемы организации учебного процесса и обосновать необходимость автоматизации;
  \item сформулировать функциональные и нефункциональные требования к системе;
  \item выполнить сравнительный анализ существующих LMS-решений и определить критерии сопоставления;
  \item обосновать выбор технологий и платформ для разработки программного комплекса;
  \item спроектировать архитектуру системы, модель данных и пользовательские сценарии;
  \item реализовать серверную, веб- и мобильную части системы;
  \item описать алгоритмы адаптивного тестирования, формирования рекомендаций и изоляции данных арендаторов;
  \item описать процедуры опытного ввода в эксплуатацию и развёртывания стенда (приложение~Г);
  \item провести испытания по программе и методике испытаний и зафиксировать ключевые результаты на критическом пути пользователей (приложение~Б);
  \item раскрыть аспекты информационной безопасности, надёжности и ограничений одного экземпляра;
  \item сформулировать выводы и подготовить приложения по ГОСТ ЕСПД (техническое задание, методика испытаний, руководство пользователя, инструкция по развёртыванию, листинги кода).
\end{parenlist}

Объектом исследования являются процессы организации корпоративного обучения, контроля прохождения курсов и оценки результатов слушателей в цифровой образовательной среде.

Предметом исследования являются методы и средства проектирования и разработки клиент--серверной информационной системы корпоративного обучения, включающей веб-интерфейс администрирования, мобильный клиент слушателя, адаптивное тестирование и механизмы аналитики.

Методы исследования: системный анализ процессов корпоративного обучения; сравнительный анализ программных аналогов по формализованным критериям и весам; объектно-ориентированное проектирование архитектуры и базы данных; документирование требований и трассировка к тест-кейсам; контроль качества по программе и методике испытаний (приложение~Б)\cite{grekul_proj_is,kulikov_testing_2024,gost7052008}.

Информационная база исследования: нормативные документы по оформлению отчётной и программной документации\cite{gost7322017,gost712003,gost21052019,gost3460190,gost250102015,gost509222006,gost1920178,gost1930179,fz152}; методические рекомендации вуза по ВКР\cite{metodichka_vkr_ikt_2021}; учебная и специальная литература по информационным системам, базам данных и тестированию\cite{fowler_patterns_ru,date_intro_db_ru,sommerville2019software,makkakov_allfusion,kaner_testing_en,van_der_linden1997,wainer2000,lord1980}; официальная документация использованных платформ\cite{fastapi_docs,postgresql_docs,flutter_docs,pydantic_docs,sqlalchemy_docs}.

Принятый вариант архитектуры --- модульный монолит на стороне сервера с раздельными клиентами и логической изоляцией данных арендаторов в общей СУБД --- обоснован в главе~2 и подтверждён опытным развёртыванием и испытаниями (приложение~Б).

Практическим результатом является программный комплекс: серверное приложение на FastAPI, веб-панель на React и TypeScript, мобильное приложение слушателя на Flutter; роли, аутентификация, управление курсами и тестами, адаптивное тестирование и рекомендации. Исходные тексты и порядок сборки приведены в открытом репозитории \url{https://github.com/lunahobi/coursum}.

Значимость результата --- возможность вести корпоративное обучение в едином контуре от контента до аналитики и повторного обучения по слабым темам.

Работа содержит введение, три главы, заключение, 55 наименований в списке источников и пять приложений по ГОСТ\cite{metodichka_vkr_ikt_2021,gost7052008}; иллюстративный материал и объём уточняются при верстке PDF.
